% ======================================================================
% SECTION 4: DISCUSSION
% Five thematic subsections: significance of the 60-second target,
% FDA framing, on-prem LLM safety advantages, head-start framing for
% the industry, and practical insights on how close the application
% is to real life today.
% Bracketed prompts for future Claude Code Opus 4.7 1M Max processing.
% Do not process the bracketed prompts here.
% ======================================================================

\subsection{Significance of the 60-Second Computational Target}
\label{sec:discussion-significance}

[DISCUSSION SUBSECTION 1 PROMPT - PROCESS DURING THE PAPER POPULATION PASS]

[Discuss the significance of the 60-second resection target as a
computational checkpoint, not as a clinical claim. The realized
quality benefit is a 240x to 480x compression of operative duration
(current operative duration 4 to 8 hours for an IDH-wildtype
glioblastoma resection; cite Stupp2005Glioblastoma, Sanai2011GBMResection),
which translates into measurable patient-experience and safety gains
under the frozen composite score formula (quality 0.40, time 0.25,
cost 0.20, safety 0.10, patient experience 0.05). Read
kevinkawchak/robotic-surgeries/2030-gbm-1min/outputs/comparison/comparison.json
and outputs/comparison_robot_vs_human/comparison.json. Render the
realized quality and safety benefits as a LaTeX table using
{>{\raggedright\arraybackslash}p{Xcm}} for each column. The columns
are: dimension, current 4 to 8 hour baseline, 60 second target,
relative improvement, source file. Highlight the per-arm
contribution analysis at outputs/viz/per_arm_contribution.png that
shows arm 1 dominates resection volume (32,400 mm cubed mean) while
arms 2 to 4 stay within the 30 percent target balance band on
coagulation seconds, suction volume, and imaging frames respectively;
this balance is the mechanical signature of the on-premises LLM
single-robot-error-minimization thesis. Cite rosa-one-brain,
Lefranc2014ROSAStereotaxy, kawchak_2026_19994945, and
repo-robotic-surgeries.]

\subsection{Respectful Yet Opportunistic FDA Framing}
\label{sec:discussion-fda}

[DISCUSSION SUBSECTION 2 PROMPT - PROCESS DURING THE PAPER POPULATION PASS]

[Discuss the FDA framing respectfully and opportunistically. The
60-second glioblastoma resection extends the FDA's 28 April 2026
Real-Time Clinical Trials proof-of-concept program (cite
fda2026realtime) from pharmacology into the surgical theater. State
explicitly that the United States needs to remain Number 1 in the
world regarding patient safety, efficacy, and speed benefits in
oncological robotic surgeries in clinical trials, and that the
on-premises LLM thesis is one route to that goal. Frame the work as
an opportunistic extension, not a critique. Reference the FDA
Software as a Medical Device (SaMD) framework (cite fda-samd) as
the appropriate regulatory pathway for an on-premises LLM that emits
control commands. Reference IEC 80601-2-77 (cite iec-80601-2-77) and
IEC 62304 (cite iec-62304) as the standards floor. Note that the
six 4-entity LLM tournament rounds at outputs/comparison/ and the
six mixed tournament rounds at outputs/comparison_robot_vs_human/
preserve the structural-time-dimension caveat in every rationale.
This caveat is the operational signature of a respectful framing: the
1-minute robot trivially beats the 1-hour human-baseline scenario on
the time dimension, but the comparison is flagged as structural and
not a fair pairwise comparison.]

\subsection{On-Premises LLM Single-Robot Error Minimization}
\label{sec:discussion-llm}

[DISCUSSION SUBSECTION 3 PROMPT - PROCESS DURING THE PAPER POPULATION PASS]

[Discuss how the on-premises LLM thesis minimizes single robot error
potential. Read
kevinkawchak/physical-ai-oncology-trials/competitions/instructions/one_minute_variant/multi_arm_coordination.md
and
kevinkawchak/robotic-surgeries/2030-gbm-1min/src/coordination/arm_heartbeat.cpp
in full. The mechanism is cross-arm safety on the deterministic 1 kHz
heartbeat broadcast bus (32-byte frames, 1 ms deadline, 3 ms
watchdog, cumulative tip force <= 12 N on the patient frame,
per-arm tip force <= 5.0 N, per-arm lateral force <= 1.0 N, 5 ms
E-stop budget, 100 microsecond emergency arm-park trigger). Render
the four-arm cross-check matrix as a LaTeX table using
{>{\raggedright\arraybackslash}p{Xcm}} for each column showing each
arm's tool, its primary sensor channels, the cross-check it performs
on its three peer arms, and the corresponding fault response. Cite
iec-80601-2-77 and iec-62304 for the regulatory floor. Reference the
on-premises LLM control loop ASCII diagram at
kevinkawchak/robotic-surgeries/2030-gbm-1min/outputs/diagrams/ that
instantiates the thesis. State explicitly that on-premises inference
keeps patient health information on site (no cloud transit) and that
the 1M token context window of Claude Code Opus 4.7 (cite
claude-opus-47, claude-code) lets a single inference pass hold all
12 instruction files, all kinematics and schema YAML / JSON files,
and all 16 iteration aggregates in context.]

\subsection{Head Start: Checking Steps the Industry No Longer Has to Take}
\label{sec:discussion-headstart}

[DISCUSSION SUBSECTION 4 PROMPT - PROCESS DURING THE PAPER POPULATION PASS]

[Discuss the head-start framing of this work. The 16-iteration
deterministic sweep, the 4-entity LLM tournament, the cross-platform
build recipes (Linux, MacOS, Windows, NVIDIA A100), and the
60-second simulation pipeline are checking steps that the industry no
longer has to take to evaluate the feasibility of an on-premises LLM
controlled 4-arm parallel stereotactic neurosurgical robot. Read
kevinkawchak/robotic-surgeries/2030-gbm-1min/outputs/reports/run_summary.md,
outputs/reports/process_log.md, and outputs/reports/final_report.md
in full. Render a LaTeX table using
{>{\raggedright\arraybackslash}p{Xcm}} for each column with columns:
checking step, current industry status, project status, source
artifact, downstream effect. Include at least the following checking
steps: 7-DOF DH parameter table per arm (industry: not published for
1-minute target; project: complete at config/kinematics_4arm.yaml),
1 kHz heartbeat broadcast bus protocol (industry: not standardized;
project: complete at src/coordination/arm_heartbeat.cpp), per-arm
force envelope and cross-check matrix (industry: per-arm only;
project: cumulative cap plus cross-check), 16-iteration deterministic
sweep with bit-determinism (industry: no published sweep at this
scale; project: complete at outputs/iterations/), structured LLM
tournament with on-premises inference (industry: cloud-only; project:
complete at outputs/comparison/ and outputs/comparison_robot_vs_human/).
Cite kawchak_2026_19994945, one-minute-variant-instructions,
repo-robotic-surgeries, kawchak_2025_17774560, kawchak_2025_15549831,
and kawchak_2025_17614396.]

\subsection{Practical Insights: How Close to Real Life Today}
\label{sec:discussion-practical}

[DISCUSSION SUBSECTION 5 PROMPT - PROCESS DURING THE PAPER POPULATION PASS]

[Discuss practical insights about how close the 60-second
glioblastoma robotic resection is to real life today. State the
distance plainly: not feasible with current commercial robotics, but
the simulation pipeline is reproducible bit for bit from seed
20260510 and runs on commodity hardware (M3 Ultra wall-clock 26 to
32 s per iteration, A100 12 s per iteration; read
2030-gbm-1min/README.md). The gaps to real life are (a) the
hypothetical NeuroSpeed 1.0 hardware does not exist (need 5x to 200x
improvements over Medtronic ROSA ONE Brain v3.0 across velocity,
acceleration, E-stop latency, positioning accuracy, and force
resolution; cite rosa-one-brain), (b) the synthetic patient PAT-GBM-0001
is not real (no IRB, no informed consent, no real EHR), and (c) the
deterministic 16-iteration sweep is not yet validated against any
real-patient cohort or TRIPOD+AI / CREMLS reporting standard (cite
Collins2024TRIPODAI, ElEmam2024CREMLS). The gaps that are close to
closing today are (i) the on-premises LLM inference backends are
production-grade for the 1M token context window (cite ollama, vllm,
claude-opus-47), (ii) the Parquet + DuckDB analytics stack is
production-grade (cite apache-arrow, duckdb), (iii) the Zenodo L0
deposition protocol is production-grade at DOI 10.5281/zenodo.18445179
(cite zenodo, repo-robotic-surgeries), and (iv) the CI gates (ruff
format, ruff check, yamllint, file size cap) are already in place
for Python 3.10, 3.11, and 3.12. Highlight that the
robotic-surgeries simulation pipeline is therefore a "ready-to-port"
software substrate: the moment a real NeuroSpeed-class robot is
available, the on-premises LLM control loop can be re-pointed at it
with minimal additional engineering.]
